% !TeX root = ../main.tex

% ==================================================
% 1.2. 파일 시스템 
% ==================================================

\subsection{파일 시스템}

\subsubsection{파일 시스템이란 무엇인가?}

컴퓨터 저장장치에 아래와 같은 파일들이 있다고 가정해보자.

\begin{dhkzsh}
사진.jpg
train.py
data.csv
requirements.txt
model.pt
\end{dhkzsh}

파일이 수십만 개가 되면 한곳에 다 넣을 수 없다.

그래서 디렉터리(directory)를 이용해서 계층적으로 정리한다.

예를 들어:

\begin{dhkdir}[in=8]
project
>data
>>train.csv
>>test.csv
>model
>>model.pt
>src
>>train.py
>>model.py
>README.md
\end{dhkdir}

이것이 우리가 앞으로 계속 보게 될 \textbf{파일 시스템의 트리 구조}다.

여기서:

\begin{dhksmry}[centered]
project \ra 디렉터리

data \ra 디렉터리

train.csv \ra 파일

src \ra 디렉터리

train.py \ra 파일
\end{dhksmry}

이며, 디렉터리 안에 다른 디렉터리가 들어갈 수 있기 때문에 트리 구조가 만들어진다.

\subsubsection{경로(Path)는 트리를 따라가는 방법이다}

위 구조에서 \texttt{train.csv}가 어디 있는지를 설명하고 싶다고 해보자.

그러면 위치를 이렇게 표현할 수 있다.

\begin{dhkzsh}
project/data/train.csv
\end{dhkzsh}

이것이 \zsh{path(경로)}다.

즉,

\begin{dhksmry}[centered]
project
\da
data
\da
train.csv
\end{dhksmry}

라는 파일 시스템의 이동 경로를 문자열 형태로 표현한 것이다.

그래서 처음에는 이렇게 생각해도 좋다.

\begin{dhknote}
\begin{dhksmry}[centered]
Path = 파일 시스템 트리에서 어떤 위치까지 찾아가는 주소
\end{dhksmry}
\end{dhknote}

단, 메모리의 주소값(address)과는 다른 개념임을 주의하자.

\subsubsection{그런데 어디서 출발하느냐가 문제가 된다}

서울에서 누군가에게

\begin{dhksmry}[centered]
"우리집은 오른쪽으로 가서 두 번째 골목이에요."
\end{dhksmry}

라고 하면 이상하다.

반면,

\begin{dhksmry}[centered]
"서울특별시 \ra ○○구 \ra ○○동 \ra ..."
\end{dhksmry}

처럼 처음부터 위치를 지정하면 출발 위치와 관계없이 찾아갈 수 있다.

파일 시스템도 똑같아서 경로에는 크게 두 종류가 있다.

\begin{dhknote}
\begin{dhksmry}[centered]
절대 경로 (absolute path)

상대 경로 (relative path)
\end{dhksmry}    
\end{dhknote}

\subsubsection{절대 경로 Absolute Path}

절대 경로는 \textbf{파일 시스템의 출발점부터 목적지까지 전부 적은 경로}다.

Mac/Linux라면 예를 들어:

\begin{dhkzsh}
/Users/user/Desktop/python/project/data/train.csv
\end{dhkzsh}

Windows 라면:

\begin{dhkzsh}
C:\Users\user\Desktop\python\project\data\train.csv
\end{dhkzsh}

이런 식이다.

Mac의 것을 구조화하면:

\begin{dhkdir}
/
>Users
>>user
>>>Desktop
>>>>python
>>>>>project
>>>>>>data
>>>>>>>train.csv
\end{dhkdir}

따라서:

\begin{dhkzsh}
/Users/user/Desktop/python/project/data/train.csv
\end{dhkzsh}

의 뜻은 다음과 같다.

\begin{dhksmry}[centered]
파일 시스템의 맨 위에서부터 train.csv까지 찾아가라.
\end{dhksmry}

\subsubsection{맨 앞의 /는 특별하다}

Mac/Linux에서 / 는 단순한 구분자가 아니다.

파일 시스템의 최상위 디렉터리인 \textbf{root directory}를 의미한다.

그래서:

\begin{dhkcode}
/Users/user/Desktop
\end{dhkcode}

를 해석하면:

\begin{dhksmry}[centered]
/
\da
Users
\da
user
\da
Desktop
\end{dhksmry}

이다.

여기서 root를 흔히 이렇게 그린다.

\begin{dhkdir}
/
>Applications
>Library
>System
>Users
>...
\end{dhkdir}

Mac을 사용하다보면 \zsh{/Users/...} 같은 경로를 자주 보게 되는 이유다.

\subsubsection{Windows는 조금 다르다}

Windows는 흔히 드라이브가 먼저 등장한다:

\begin{dhkcode}
C:\Users\user\Desktop
\end{dhkcode}

\subsubsection{상대 경로 Relative Path}

현재 내가 있는 위치가:

\begin{dhkcode}
project/
\end{dhkcode}

라고 하자.

구조는:

\begin{dhkdir}
project
>data
>>train.csv
>src
>>train.py
\end{dhkdir}

현재 위치가 \zsh{project}라면 \zsh{train.csv}를 굳이:

\begin{dhkcode}
/Users/user/Desktop/python/project/data/train.csv
\end{dhkcode}

라고 전부 적을 필요가 없이:

\begin{dhkcode}
data/train.csv
\end{dhkcode}

라고 적으면 된다.

왜냐하면 출발점이 이미 \zsh{project}이기 때문이다.

이것이 \textbf{상대 경로}다.

\subsubsection{그러면 "현재 위치"가 대체 무엇인가?}

여기가 Python에서 굉장히 중요하다.

이것을:

\begin{dhknote}
\begin{dhksmry}[centered]
Current Working Directory, CWD
\end{dhksmry}
\end{dhknote}

라고 한다.

한국어로는 \textbf{현재 작업 디렉터리}정도로 번역이 된다.

예를 들어 터미널이:

\begin{dhkzsh}
~/Desktop/python/project %
\end{dhkzsh}

라고 되어 있다면 현재 작업 디렉터리가 대략:

\begin{dhkzsh}
/Users/user/Desktop/python/project
\end{dhkzsh}

라고 생각할 수 있다.

Python에서도 확인할 수 있다.

\begin{dhkcode}
from pathlib import Path

print(Path.cwd())
\end{dhkcode}

예를 들어:

\begin{dhkzsh}
/Users/user/Desktop/python/project
\end{dhkzsh}

가 출력된다.

\subsubsection{상대 경로는 CWD를 기준으로 해석한다}

이 원칙이 정말 중요하다.

\begin{dhkcode}
from pathlib import Path

p = Path("data/train.csv")
\end{dhkcode}

여기서:

\begin{dhkzsh}
data/train.csv
\end{dhkzsh}

은 상대 경로다.

만약:

\begin{dhkcode}
Path.cwd()
\end{dhkcode}

가:

\begin{dhkzsh}
/users/user/Desktop/python/project
\end{dhkzsh}

라면 실질적으로 가리키는 곳은:

\begin{dhksmry}[centered]
/users/user/Desktop/python/project
+
data/train.csv
\end{dhksmry}

즉:

\begin{dhkzsh}
/users/user/Desktop/python/project/data/train.csv
\end{dhkzsh}

이다.

\begin{dhkdir}
/Users/user/Desktop/python/@project (CWD)
>data
>>train.csv
\end{dhkdir}

이 원리를 잡으면 나중에 Pandas에서 발생하는 수많은 파일 경로 오류를 이해할 수 있다.

\subsubsection{. 은 "현재 디렉터리"다}

경로에서 이것도 많이 본다.

\begin{dhkzsh}
.
\end{dhkzsh}

점 하나는:

\begin{dhksmry}[centered]
현재 디렉터리
\end{dhksmry}

를 의미한다.

따라서:

\begin{dhkzsh}
./data/train.csv
\end{dhkzsh}

는:

\begin{dhksmry}[centered]
현재 디렉터리
\da
data
\da
train.csv
\end{dhksmry}

라는 뜻이다.

그래서:

\begin{dhkzsh}
data/train.csv
\end{dhkzsh}

와

\begin{dhkzsh}
./data/train.csv
\end{dhkzsh}

는 이 상황에서는 사실상 같은 위치를 나타낸다.

\subsubsection{.. 는 부모 디렉터리다}

점 두 개:

\begin{dhkzsh}
..
\end{dhkzsh}

는:

\begin{dhksmry}[centered]
현재 디렉터리의 한 단계 위
\end{dhksmry}

를 의미한다.

예를 들어 위치가:

\begin{dhkzsh}
project/src
\end{dhkzsh}

라고 하자.

구조는:

\begin{dhkdir}
project
>data
>>train.csv
>src
>>train.py
\end{dhkdir}

현재 위치:

\begin{dhkzsh}
project/src
\end{dhkzsh}

에서:

\begin{dhkzsh}
..
\end{dhkzsh}

하면 \zsh{project}로 올라간다.

그러므로:

\begin{dhkzsh}
../data/train.csv
\end{dhkzsh}

는:

\begin{dhkbranch}[in=14]
node main :: src
arrow main :: ..
node main :: project
arrow main :: data
node main :: data
arrow main
node main :: train.csv
\end{dhkbranch}

는 뜻이다.

즉,

\begin{dhksmry}[centered]
src \ra 한 단계 위 \ra data \ra train.csv
\end{dhksmry}

\subsubsection{여러 번 올라갈 수도 있다}

예를 들어:

\begin{dhkdir}
project/
>src/
>>models/
>>>transformer/
\end{dhkdir}

현재 위치가:

\begin{dhkzsh}
project/src/models/transformer
\end{dhkzsh}

라고 하자.

그러면:

\begin{dhkzsh}
..
\end{dhkzsh}

는:

\begin{dhkzsh}
project/src/models
\end{dhkzsh}

이고,

\begin{dhkzsh}
../..
\end{dhkzsh}

는:

\begin{dhkzsh}
project/src
\end{dhkzsh}

이고,

\begin{dhkzsh}
../../..
\end{dhkzsh}

는:

\begin{dhkzsh}
project
\end{dhkzsh}

이다.

따라서:

\begin{dhkzsh}
../../../data/train.csv
\end{dhkzsh}

같은 것도 가능하다.

하지만 실제 프로젝트에서 이런 경로가 난무하기 시작하면 구조를 다시 생각해볼 필요가 있다.

\subsubsection{~ 는 또 다른 의미다}

Mac/Linux 터미널에서 자주 보는:

\begin{dhkzsh}
~ 
\end{dhkzsh}

는 보통 사용자의 \textbf{home directory}를 의미한다.

Mac 계정 이름을 예시로 dhk 라고 하면:

\begin{dhkzsh}
~
\end{dhkzsh}

가 대략:

\begin{dhkzsh}
/Users/dhk
\end{dhkzsh}

를 의미한다.

그래서:

\begin{dhkzsh}
~/Desktop
\end{dhkzsh}

은:

\begin{dhkzsh}
/Users/dhk/Desktop
\end{dhkzsh}

같은 의미가 된다.

Windows에도 사용자 홈 개념이 있고 Python의 pathlib으로 운영체제에 맞게 처리할 수 있다.

예를 들어:

\begin{dhkcode}
from pathlib import Path
print(Path.home())
\end{dhkcode}

이라고 하면 사용자 홈 디렉터리를 얻을 수 있다.

\subsubsection{여기서 중요한 함정}

다음 파일이 있다고 해보자.
\begin{dhkdir}
project
>data
>>train.csv
>src
>>train.py
\end{dhkdir}

\zsh{train.py}에:

\begin{dhkcode}
import pandas as pd

df = pd.read_csv("../data/train.csv")
\end{dhkcode}

라고 작성했다.

많은 사람이 이렇게 생각한다.
\begin{dhksmry}[centered]
train.py가 src안에 있으니까 .. 하면 project겠네.
\end{dhksmry}

\textbf{항상 그런 것은 아니다}.

상대 경로의 기본 기준은 \textbf{Python 파일이 위치한 디렉터리가 아니라 현재 작업 디렉터리(CWD)다}.

이게 처음 파일 경로를 배울 때 가장 많이 혼동하는 부분 중 하나다.

\subsubsection{실행하는 위치가 중요하다}

예를 들어 터미널에서:

\begin{dhkzsh}
project/src %
\end{dhkzsh}

인 상태에서:

\begin{dhkzsh}
python train.py
\end{dhkzsh}

를 실행했다고 하자.

CWD:

\begin{dhkzsh}
project/src
\end{dhkzsh}

따라서:

\begin{dhkzsh}
../data/train.csv
\end{dhkzsh}

는:
\begin{dhkbranch}[in=10]
node main :: project/src
arrow main :: ..
node main :: project
arrow main :: data
node main :: project/data/train.csv
\end{dhkbranch}


가 되며 잘 작동한다.

그런데 이번에는:

\begin{dhkzsh}
project %
\end{dhkzsh}

에서:

\begin{dhkzsh}
python src/train.py
\end{dhkzsh}

를 실행했다고 해보자.

\zsh{train.py} 파일은 여전히:

\begin{dhkzsh}
project/src/train.py
\end{dhkzsh}

에 있다.

하지만 CWD는 \zsh{project}다.

따라서:

\begin{dhkzsh}
../data/train.csv
\end{dhkzsh}

를 해석하면:

\begin{dhksmry}[centered]
project
\da
project의 부모
\da
data/train.csv
\end{dhksmry}

를 찾게 된다.

\subsubsection{이것 때문에 파일은 분명히 있는데 못찾는 경우가 발생한다}

아마 Python이나 ML을 하면서 \zsh{FileNotFoundError}를 본적이 있을 것이다.

파일이 정말 없는 경우도 있지만 상당히 흔한 상황은:

\begin{dhknote}
\begin{dhksmry}[centered]
내가 생각한 기준 디렉터리 $\neq$ Python의 현재 작업 디렉터리
\end{dhksmry}
\end{dhknote}

인 것이다.

그래서 파일을 못 찾을 때 가장 먼저:

\begin{dhkcode}
from pathlib import Path

print(Path.cwd())
\end{dhkcode}

를 찍어보는 습관이 상당히 유용하다.

\subsubsection{이제 \zsh{pathlib.Path}를 사용해보자}

예전 Python 코드에서는:

\begin{dhkcode}
import os
\end{dhkcode}

와 문자열 경로 조작이 굉장히 많았다.

현재는 경로 자체를 다룰 때 \zsh{pathlib}이 매우 편하다.

\begin{dhkcode}
from pathlib import Path
\end{dhkcode}

그리고:

\begin{dhkcode}
p = Path("data/train.csv")
\end{dhkcode}

라고 할 수 있다.

여기서 중요한 사고방식이 있다.

\begin{dhkzsh}
"data/train.csv"
\end{dhkzsh}

는 그냥 문자열이다.

반면:

\begin{dhkcode}
Path("data/train.csv")
\end{dhkcode}

는 경로를 표현하고 조작하기 위한 Path 객체다.

\subsubsection{문자열과 \zsh{Path} 객체의 차이}

문자열:

\begin{dhkcode}
p = "data/train.csv"
\end{dhkcode}

Python 입장에서는 그냥 문자들의 나열이다.

\begin{dhkzsh}
    d a t a / t r a i n . c s v
\end{dhkzsh}

그 이상을 특별히 알 필요가 없다.

반면:

\begin{dhkcode}
p = Path("data/train.csv")
\end{dhkcode}

라고 하면 경로 관련 기능을 사용할 수 있다.

예를 들어:

\begin{dhkcode}
p.exist()
\end{dhkcode}

파일이 실제 존재하는지 확인한다.

\begin{dhkcode}
p.is_file()
\end{dhkcode}

파일인지 확인한다.

\begin{dhkcode}
p.is_dir()
\end{dhkcode}

디렉터리인지 확인한다.

\begin{dhkcode}
p.name
\end{dhkcode}

이라면:

\begin{dhkzsh}
train.csv
\end{dhkzsh}

를 얻는다.

\subsubsection{경로를 / 로 조립할 수도 있다}

\zsh{pathlib}의 편한 기능이다.

\begin{dhkcode}
from pathlib import Path

project = Path("/Users/user/Desktop/project")

data = project / "data"
train = data / "train.csv"
\end{dhkcode}

결과적으로:

\begin{dhkzsh}
/Users/user/Desktop/project/data/train.csv
\end{dhkzsh}

라는 경로가 만들어진다.

처음 보면:

\begin{dhkcode}
project / "data"
\end{dhkcode}

가 이상할 수 있다.

숫자 나눗셈이 아니다.

\begin{dhknote}
\begin{dhksmry}[centered]
\zsh{Path}객체가 / 연산자를 \textbf{경로 연결이라는 의미로 정의해 놓은 것}이다.
\end{dhksmry}
\end{dhknote}

그래서:

\begin{dhkcode}
Path("data") / "train.csv"
\end{dhkcode}

는:

\begin{dhkzsh}
data/train.csv
\end{dhkzsh}

라는 경로가 된다.

\subsubsection{절대 경로인지 확인할 수도 있다}

\begin{dhkcode}
from pathlib import Path

a = Path("/Users/user/Desktop/data.csv")
b = Path("data/data.csv")

print(a.is_absolute())
print(b.is_absolute())
\end{dhkcode}

결과:

\begin{dhkzsh}
True
False
\end{dhkzsh}

왜냐하면:

\begin{dhksmry}
/Users/user/Desktop/data.csv
\ua
\zsh{root}부터 시작
\end{dhksmry}

하지만:

\begin{dhkzsh}
data/data.csv
\end{dhkzsh}

는 출발점이 없다.

따라서 \zsh{CWD}가 필요하다.

\subsubsection{상대 경로를 절대 경로로 확인해보자}

\begin{dhkcode}
from pathlib import Path

p = Path("data/train.csv")

print(p)
print(p.resolve())
\end{dhkcode}

예를 들어 \zsh{CWD}가:

\begin{dhkzsh}
/Users/user/Desktop/project
\end{dhkzsh}

라면:

\begin{dhkcode}
print(p)
\end{dhkcode}

는:

\begin{dhkzsh}
data/train.csv
\end{dhkzsh}

이고,

\begin{dhkcode}
print(p.resolve())
\end{dhkcode}

는 대략:

\begin{dhkzsh}
/Users/user/Desktop/project/data/train.csv
\end{dhkzsh}

가 된다.

즉 \zsh{resolve()}를 사용하면:
\begin{dhksmry}[centered]
"이 상대 경로가 결국 실제로 어디를 가리키는 거야?"
\end{dhksmry}

를 확인하기 좋다.

\subsubsection{실제 ML 프로젝트를 하나 만들어보자}

앞으로 이 예제를 확장해보자

\begin{dhkdir}[in=10]
my_ml_proect/
>data/
>>train.csv
>>test.csv
>models/
>notebooks/
>>experiment.ipynb
>src/
>>train.py
>>utils.py
>README.md
\end{dhkdir}

터미널에서:

\begin{dhkzsh}
my_ml_project %
\end{dhkzsh}

에 있다고 하자.

즉, \zsh{CWD}:

\begin{dhkzsh}
my_ml_project
\end{dhkzsh}

이다.

그러면 Python에서:

\begin{dhkcode}
from pathlib import Path

cwd = Path.cwd()
print(cwd)
\end{dhkcode}

대략:

\begin{dhkzsh}
/Users/user/.../my_ml_project
\end{dhkzsh}

이다.

그리고:

\begin{dhkcode}
data_dir = cwd / "data"
train_file = data_dir / "train.csv"
\end{dhkcode}

그러면:

\begin{dhkzsh}
my_ml_project/data/train.csv
\end{dhkzsh}

를 가리키게 된다.

\subsubsection{Pandas와도 바로 연결된다}

\zsh{Path}객체는 Pandas같은 현대 Python 라이브러리에 바로 사용할 수 있다.

\begin{dhkcode}
from pathlib import Path
import pandas as pd

data_path = Path("data") / "train.csv"

df = pd.read_csv(data_path)
\end{dhkcode}

이제 코드의 의미를 완전히 분해할 수 있다.

\begin{dhkcode}
Path("data")
\end{dhkcode}

현재 작업 디렉터리를 기준으로 \zsh{data}라는 경로를 나타내고,

\begin{dhkcode}
/ "train.csv"
\end{dhkcode}

를 통해 그 아래의 파일 경로를 만들고,

\begin{dhkcode}
pd.read_csv(data_path)
\end{dhkcode}

가 해당 위치의 csv를 읽는다.

\subsubsection{디렉터리도 만들 수 있다}

예를 들어 학습된 모델을:

\begin{dhkzsh}
models/
\end{dhkzsh}

에 저장하고 싶다고 하자.

\begin{dhkcode}
from pathlib import Path

model_dir = Path("models")
model_dir.mkdir()
\end{dhkcode}

그러면 디렉터리를 만들 수 있다.

이미 존재할 가능성이 있다면:

\begin{dhkcode}
model_dir.mkdir(exist_ok=True)
\end{dhkcode}

를 사용할 수 있다.

더 깊은 구조:

\begin{dhkzsh}
outputs/models/checkpoints
\end{dhkzsh}

를 한 번에 만들고 싶다면:

\begin{dhkcode}
path = Path("outputs/models/checkpoints")
path.mkdir(parens=True, exist_ok=True)
\end{dhkcode}

같은 식으로 사용할 수 있다.

이런 패턴은 ML/DL 프로젝트에서 굉장히 흔하다.

\subsubsection{이제 파일 경로와 URL을 절대로 섞지 말자}

이것도 앞으로 매우 중요하다.

다음은 \textbf{로컬 파일 경로}다.

\begin{dhkzsh}
/users/user/project/data/train.csv
\end{dhkzsh}

또는:

\begin{dhkzsh}
C:\Users\user\project\data\train.csv
\end{dhkzsh}

컴퓨터의 파일 시스템에서 위치를 찾는다.

반면:

\begin{dhkzsh}
https://example.com/data/train.csv
\end{dhkzsh}

는 URL이다.

이것은:

\begin{dhksmry}[centered]
내 SSD
\da
디렉터리 탐색
\end{dhksmry}

의 문제가 아니다.

대략적으로:

\begin{dhksmry}[centered]
URL
\da
네트워트
\da
서버
\da
HTTP 요청
\da
응답
\end{dhksmry}

의 세계다.

따라서 나중에 우리가 \zsh{Path}와 \zsh{URL}을 명확히 분리해서 배울 것이다.

\subsubsection{지금까지의 개념을 하나의 그림으로 합치자}

파일 시스템이:

\begin{dhkdir}
/
>Users
>>user
>>>Desktop
>>>>my_ml_project
>>>>>data
>>>>>>train.csv
>>>>>models
>>>>>src
>>>>>>train.py
\end{dhkdir}

라고 하자.

현재 작업 디렉터리:

\begin{dhkzsh}
/Users/user/Desktop/my_ml_project
\end{dhkzsh}

이라면:

\begin{dhkzsh}
절대 경로
/Users/user/Desktop/my_ml_project/data/train.csv
\end{dhkzsh}

와:

\begin{dhkzsh}
상대 경로
data/train.csv
\end{dhkzsh}

는 같은 파일을 가리킬 수 있다.

그리고:

\begin{dhkzsh}
.
\end{dhkzsh}

는:

\begin{dhkzsh}
/Users/user/Desktop/my_ml_project
\end{dhkzsh}

를 의미하므로:

\begin{dhkzsh}
./data/train.csv
\end{dhkzsh}

도 같은 파일이다.

반면:

\begin{dhkzsh}
..
\end{dhkzsh}

는:

\begin{dhkzsh}
/Users/user/Desktop
\end{dhkzsh}

을 의미한다.

\subsubsection{여기까지에서 가장 중요한 사고방식}

위 내용들을 단순히 기호로 외우면 안된다.

다음 구조를 머릿속에 만드는 것이 목표다.

\begin{dhksmry}[centered]
컴퓨터 저장장치
\da
파일 시스템
\da
디렉터리의 계층적 트리
\da
트리상의 위치를 표현
\da
Path
\end{dhksmry}

\begin{dhkdir}
Path
> 상대 경로 - CWD 기준 . / ..
> 절대 경로 - root 기준
\end{dhkdir}

그리고 Python에서는:

\begin{dhksmry}[centered]
Path.cwd()
\da
현재 작업 디렉터리

Path("data/train.csv")
\da
상대 경로

Path("/Users/.../train.csv")
\da
절대 경로

.resolve()
\da
최종적으로 어디를 가리키는지 확인
\end{dhksmry}

이렇게 연결된다.

\subsubsection{다음 단계가 특히 중요하다}

바로 \zsh{sys}로 가지 않고 \textbf{경로에서 가장 많이 혼동하는 세 가지 기준점}을 분리하는 게 좋다.

\begin{dhknote}
\begin{dhksmry}
1. 현재 작업 디렉터리 \ra Path.cwd()

2. 실행 중인 .py 파일의 위치 \ra __file__

3. Python이 import 대상을 찾는 위치 \ra sys.path
\end{dhksmry}
\end{dhknote}

셋 다 화면에는 \zsh{/Users/...} 같은 경로가 나오기 때문에 같은 종류처럼 보이지만, 실제 의미는 완전히 다르다.

여기까지 분리되면 그 다음 \zsh{sys.argv}를 배울 때

\begin{dhkzsh}
python src/train.py --data data/trian.csv
\end{dhkzsh}

에서 \textbf{shell $\rightarrow$ Python interpreter $\rightarrow$ .py 파일 $\rightarrow$ Path $\rightarrow$ 실제 파일}이 한번에 연결되기 시작한다.



